\section{Descripción del problema y modelo}

\subsection{Grilla, viviendas y hogares}

Montevideo se representa mediante una grilla bidimensional rectangular con celdas residenciales y no residenciales, siguiendo la abstracción espacial del modelo clásico de Schelling. Una celda residencial contiene como máximo una vivienda, que puede estar ocupada o vacía. Los hogares se clasifican en siete subestratos y tres clases socioeconómicas.

El escenario completo inicial utiliza $1024\times640$ celdas. Sus coordenadas expresan proximidad discreta dentro del modelo y no una proyección cartográfica ni ubicaciones reales de Montevideo. Los ocho municipios se representan como zonas contiguas dimensionadas a partir de datos agregados. Esta simplificación permite conservar totales oficiales sin afirmar que la grilla reproduce la forma, las distancias o la trama urbana de la ciudad. La grilla utiliza bordes cerrados: los vecindarios se recortan en los límites y los extremos opuestos no se conectan.

\subsection{Vecindario y ponderación}

Para una celda residencial $c$, el vecindario de radio $r$ se define como
\begin{equation}
N_r(c)=\{c'\in C_R : d(c,c')\leq r\}.
\end{equation}
La influencia disminuye con la distancia mediante un kernel gaussiano,
\begin{equation}
w(c,c')=\exp\left(-\frac{d(c,c')^2}{2\sigma^2}\right).
\end{equation}

El escenario base utiliza pertenencia de Chebyshev con $r=2$, por lo que considera un cuadrado de hasta 24 vecinos, y calcula el peso con distancia euclidiana y $\sigma=1{,}0$. Separar ambas métricas permite un recorrido rectangular eficiente y, a la vez, asignar menor influencia a diagonales y celdas lejanas. Como referencia se evaluará una variante clásica con $r=1$ y ponderación uniforme.

\subsection{Satisfacción y tolerancias}

La proporción ponderada de cada clase se compara con una matriz configurable de preferencias mínimas y tolerancias máximas. La propia unidad familiar no participa del cálculo.

Cuando no existe ningún hogar vecino, la unidad familiar se considera satisfecha por aislamiento. La misma regla se aplica al evaluar una vivienda candidata y estos casos se contabilizan por separado en las métricas. Esta convención evita asignar proporciones artificiales a un denominador vacío.

La configuración base exige para una unidad de clase alta al menos 70\% de vecinos altos, como máximo 30\% medios y 5\% bajos; para una unidad media, como máximo 70\% altos, al menos 50\% medios y como máximo 40\% bajos; y para una unidad baja, como máximo 80\% altos, 90\% medios y al menos 2\% bajos. Todas las condiciones de la fila deben cumplirse simultáneamente.

\subsection{Precios y restricción económica}

El precio propuesto sigue un modelo hedónico simplificado:
\begin{equation}
\log p_v(t)=\alpha_0+\beta_1D_v(t)+\beta_2A_v(t)+\epsilon_v(t).
\end{equation}
La demanda $D_v(t)$ es la proporción ponderada de viviendas residenciales ocupadas del vecindario. El poder adquisitivo $A_v(t)$ es el ingreso medio ponderado de sus hogares, normalizado entre los ingresos de referencia de B$-$ y A+. Ambas variables pertenecen a $[0,1]$; cuando no existe un denominador válido se asigna cero. El ruido $\epsilon_v(t)$ sigue una normal con media cero y desviación $0{,}05$, truncada a $[-0{,}15;0{,}15]$, y se genera de forma determinista por vivienda e iteración mediante un generador sin estado compartido.

El precio y los ingresos se expresan en miles de pesos constantes de 2022. Para aproximar la cuota se usa una anualidad de cuota fija sobre el 80\% del precio, con tasa nominal anual de 6\% y plazo de 240 meses. Estos parámetros son configurables y no pretenden reproducir un producto hipotecario real. Una vivienda es accesible cuando su cuota no supera la fracción $\rho$ del ingreso del hogar.

La configuración base adopta $\beta_1=0{,}40$, $\beta_2=0{,}60$ y $\rho=0{,}30$. El intercepto no se selecciona arbitrariamente: se calibra para que un hogar M, en un entorno con $D=0{,}90$ y el poder adquisitivo normalizado correspondiente a su ingreso, pague una cuota equivalente al 25\% de dicho ingreso. Con los supuestos financieros anteriores se obtiene un precio de referencia de $1448{,}15$ miles de pesos de 2022 y $\alpha_0=6{,}793886203$. Esta cuota objetivo deja un margen respecto del límite de accesibilidad del 30\%.

\subsection{Mudanzas y evolución temporal}

La simulación es síncrona. Los hogares insatisfechos y no bloqueados solicitan destinos accesibles; los conflictos se resuelven de forma determinista antes de aplicar simultáneamente las mudanzas. Una vivienda liberada queda disponible en la iteración siguiente.

Después de mudarse en el mes $t$, un hogar permanece bloqueado durante los 12 meses siguientes y vuelve a ser elegible en $t+13$. Su satisfacción se continúa evaluando durante ese período. El parámetro es configurable y se considera una variante sin permanencia mínima.
